\section{Conclusion}

In this project, EuroSAT RGB is used for two related classification tasks. The
first completed model is a CNN implemented and trained from scratch. Its final
configuration was not chosen arbitrarily, but selected from several validation
experiments. SGD with momentum, combined with augmentation, dropout, and weight
decay, gave the best validation results for the scratch CNN.

The final scratch CNN achieved a test accuracy of 0.9531 and a macro-F1 of
0.9519 on the 10-class task. On the binary task, it achieved 0.9795 accuracy and
0.9795 macro-F1. The results show that a compact, self-trained CNN works well on
EuroSAT RGB. The remaining errors mainly occur between visually similar classes,
especially within the vegetation categories.

Model 2 is a transfer learning approach, evaluated with two ImageNet-pretrained
backbones. The first, EfficientNet-B0, had its classifier replaced for the
EuroSAT tasks and was trained in two phases: first with a frozen feature
extractor and then with the final EfficientNet blocks unfrozen.

The second backbone for Model 2, ResNet-50, was evaluated to test whether the
transfer-learning advantage holds across architectures. It achieved a test
accuracy of 0.9817 and a macro-F1 of 0.9812 on the
10-class task, and 0.9938 on the binary task. The data efficiency experiments
show that the pretrained model reaches competitive accuracy with far less
labeled data: at only 500 training images it achieves 92.8\% test accuracy on
the 10-class task, compared to 70.2\% for the scratch CNN. The gap narrows when
the full dataset is available, but the ResNet consistently outperforms the
scratch CNN across all data sizes and both tasks.

The results show that both approaches work well on EuroSAT. Both Model 2 backbones
outperformed the scratch CNN on both tasks. ResNet-50 achieved the best 10-class
result (macro-F1 0.9812), while EfficientNet-B0 was marginally best on the binary
task (macro-F1 0.9941, versus 0.9938 for ResNet-50). The binary task was easier
for all models because
several visually similar
vegetation classes were merged into one class. The remaining errors mainly occur
between visually similar land use categories.

A useful extension would be to use all Sentinel-2 spectral bands instead of RGB
only. This would require either adapting the first layer of a pretrained model
or using a model designed for multispectral satellite images.